% 第 6 章 Reports（原书 6-102 … 6-121）
\bichapter{报告}{Reports}
\label{chap:reports}

\section{引言}
\subsection*{Introduction}

完成所选分析后，有多种方法可查看结果。Tools 与 Results 菜单提供多种结果色阶图（colormap）显示方式。详见 GUI 菜单说明（附录 D「RedHawk Graphic User Interface Description」）。

RedHawk 将每条分析命令的进度与结果输出到标准文本流。此外，摘要报告写入 \texttt{adsRpt} 与 \texttt{adsPower} 目录。本节描述各类 RedHawk 结果、地图与报告文件。

最完整的结果与输出分析工具集位于 Explorer 实用程序（Explorer 菜单下）。

\section{RedHawk 日志文件}
\subsection*{RedHawk Log Files}

\subsection{命令文件}
\subsubsection*{Command Files}

每个 RedHawk 会话在 \texttt{adsRpt/Cmd} 目录下生成命令文件，记录所有已执行的一级命令，例如设计设置、提取、功耗计算与静态 IR drop 分析。文件名中的日期与时间戳扩展名表示 RedHawk 会话启动时刻：
\begin{lstlisting}
redhawk.cmd.<date>-<time-stamp>
\end{lstlisting}
若需回放，可将该命令文件用作批处理回放文件。

\subsection{日志与错误文件}
\subsubsection*{Log and Error Files}

每个 RedHawk 会话在 \texttt{adsRpt/Log} 目录下生成完整日志文件，文件名扩展名为会话启动日期与时间，例如：
\begin{lstlisting}
redhawk.log.<date>-<start-time>
\end{lstlisting}
日志文件记录完整 RedHawk 运行的标准输出，包括全部 Error 与 Warning 消息。会话信息亦显示在日志消息窗口并写入日志文件。

Error 与 Warning 消息还分别写入 \texttt{adsRpt/Error} 与 \texttt{adsRpt/Warn} 子目录下的独立文件：
\begin{lstlisting}
Warn/redhawk.warn.<date>-<time-stamp>
Error/redhawk.err.<date>-<time-stamp>
\end{lstlisting}
可使用 \texttt{-log} 选项将 RedHawk 日志与 error/warning 文件重定向到其他目录或文件名：
\begin{lstlisting}
redhawk -log <user_log_dir/filename>
\end{lstlisting}
这会在所选目录与文件名下创建日志文件，并在同一目录下创建名称类似（无时间戳）的 Warning 与 Error 文件：
\begin{itemize}
  \item Log：\texttt{<user\_log\_dir/filename>}
  \item Warning：\texttt{<user\_log\_dir/filename>.warn}
  \item Error：\texttt{<user\_log\_dir/filename>.err}
\end{itemize}

\subsection{最新日志与消息文件的链接}
\subsubsection*{Links to Latest Log and Message Files}

在 \texttt{adsRpt/} 目录下，RedHawk（以及 GDS2DEF/GDS2RH、APL 等其他工具）创建顶层 Log、Command、Warning 与 Error 文件，作为最新版本文件的软链接：
\begin{itemize}
  \item Log：\texttt{Log/redhawk.log.<date>-<time-stamp>}（链接至最新文件）
  \item Command：\texttt{Cmd/redhawk.cmd.<date>-<time-stamp>}
  \item Error：\texttt{Error/redhawk.err.<date>-<time-stamp>}
  \item Warning：\texttt{Warn/redhawk.warn.<date>-<time-stamp>}
\end{itemize}
APL、gds2def/gds2rh 及其他 RedHawk 工具遵循相同协议。

\subsection{GUI 日志消息查看器}
\subsubsection*{GUI Log Message Viewer}

在 GUI 中，可通过 Results 菜单选择 Log Message Viewer，快速查找特定 Info、Warning 与 Error 报告，如图~\ref{fig:rh-log-viewer-open} 所示。

\figplaceholder{Figure 6-6}{打开 Log Message Viewer}
\label{fig:rh-log-viewer-open}

默认情况下，上窗口显示当前会话生成的 Info、Warning 与 Error 日志消息摘要，如图~\ref{fig:rh-log-messages} 所示。

点击某一 Log Message 类别可列出该类别下的消息（图~\ref{fig:rh-error-list}）。再点击单条消息时，该消息在下窗口的会话日志中高亮。在上窗口对消息右键可获取该 Error 或 Warning 的附加信息。

点击 CPU/Memory Usage 可查看截至当前各阶段的会话 CPU 与内存使用摘要。

Setup Design 按钮显示设置步骤中处理的所有输入文件，并以颜色编码标示处理是否有问题：
\begin{itemize}
  \item 绿色——处理该文件时未遇到问题
  \item 橙色——处理时产生一条或多条 Warning
  \item 红色——处理时产生一条或多条 Error
\end{itemize}

\figplaceholder{Figure 6-7}{Log Viewer 中的消息列表}
\label{fig:rh-log-messages}

Power 按钮显示功耗计算步骤的 Power Summary 报告。

Log Message Viewer 底部 Log File 窗口显示当前所展示日志文件的路径。Browse 与 Apply 按钮可选择并显示其他 RedHawk 会话的日志。若会话在显示日志后继续运行，可用 Refresh 按钮更新日志消息显示。

\figplaceholder{Figure 6-8}{查看 Error 消息列表}
\label{fig:rh-error-list}

\subsection{审查设计中的短路}
\subsubsection*{Reviewing Shorts in the Design}

RedHawk 可检测并报告任意两条网络之间的短路，包括同域（例如均为电源网络）或不同域（一电源一地）情形。短路可发生在导线之间、导线与 via、导线与 pin、via 与 pin 等。在 GUI 中可通过 View $\rightarrow$ Connectivity $\rightarrow$ Shorts 显示短路，列出设计中所有短路，点击列表项可缩放到对应位置。

RedHawk 亦在 \texttt{adsRpt/Error/redhawk.err<timestamp>} 文件中列出短路，消息 ID 为 CON-109、CON-110 与 CON-111。报告情形取决于被短路对象是否均被提取，以及 GSR 中 \gsr{IGNORE_SHORT} 关键字的设置。

\begin{enumerate}
  \item 若一条被短路网络需被提取（在 GSR 的 \gsr{VDD_NET} 或 \gsr{GND_NET} 列表中），另一条不需提取（不在 \gsr{VDD_NET} 与 \gsr{GND_NET} 列表中），则无论 \gsr{IGNORE_SHORT} 如何设置，RedHawk 均报告 CON-111：
\begin{lstlisting}
WARNING (CON-111): A short between net <A> and <B> is detected on layer
<M1> at (x1,y1 x2,y2) by wires. But since net<B> is not extracted,
net<A>'s R network is not affected.
\end{lstlisting}
  \item 若两条被短路网络均被提取（在 GSR 的 \gsr{VDD_NET} 或 \gsr{GND_NET} 列表中），则根据 GSR 中 \gsr{IGNORE_SHORT} 设置报告 CON-109 或 CON-110：
  \begin{enumerate}
    \item \gsr{IGNORE_SHORT} 1（默认）：RedHawk \textbf{不}将两条网络的 R 网络短路，并报告 CON-110：
\begin{lstlisting}
WARNING (CON-110): A short between net <A> and <B> is detected on
layer<M1> at (x1,y1 x2,y2) by wires. The short is ignored and their
R network will not be shorted!
\end{lstlisting}
    \item \gsr{IGNORE_SHORT} 0：RedHawk 将两条网络的 R 网络短路，并报告 CON-109：
\begin{lstlisting}
ERROR (CON-109): A short between net<A> and <B> is detected on
layer<M1> at (x1,y1 x2,y2) by wires, their R network will be shorted!
\end{lstlisting}
  \end{enumerate}
\end{enumerate}

\subsection{Tech 文件查看器}
\subsubsection*{Tech File Viewer}

通过 GUI 的 Tech File Viewer 可呈现 Tech 文件的不同视图。使用 View $>$ Technology Layers 查看 Tech 文件中层参数信息（图~\ref{fig:rh-tech-viewer}）。Technology Viewer 提供 Zoom 按钮放大信息，以及 GIF 按钮导出 Technology Viewer 的 GIF 格式输出。

\figplaceholder{Figure 6-9}{Tech File Viewer 用法}
\label{fig:rh-tech-viewer}

还可在 technology viewer 对话框中通过「Substrate button」激活，在 technology viewer 中查看两层金属之间最多三种 via 模型。

\subsection{DEF 导入摘要}
\subsubsection*{DEF Import Summary}

导入 DEF 后创建 \texttt{adsRpt/tech\_summary.rpt}，报告设计中使用的层与被忽略的层。

\section{功耗摘要文件}
\subsection*{Summary Files for Power}

\subsection{\texttt{power\_summary.rpt} 文件}
\subsubsection*{\texttt{power\_summary.rpt} File}

GUI 中 Power $>$ Calculate Power，或 TCL shell 中 \cmd{perform pwrcalc}，会自动在 \texttt{adsRpt} 目录生成功耗摘要文件 \texttt{power\_summary.rpt}，包含芯片各 block 的详细功耗分解。示例如下：
\begin{lstlisting}
-------------------------------------------------------------------
Recommended dynamic simulation time, 5000psec, to include 100%
of total power for DYNAMIC_SIMULATION_TIME in GSR.

INFO: Importing user specified power file: demo.inst.power
INFO: Instances not specified in this file will have zero power.
INFO: 100% of instances specified in instance_power_file have
corresponding instances in the design.
0 out of 4142 instances do not have corresponding instance in
the design.
INFO: For complete list of PWR-121 and PWR-122 WARNINGS, please
see adsRpt/apache.inst_pwr_file.mismatch.

The power is based on the INSTANCE_POWER FILE: demo.inst.power
Redhawk honors instance-by-instance power as specified in the
above file.
Individual components of power, like switching power, internal
power, however, are calculated by Redhawk and may have been
scaled to meet the total instance-specific power number.

Power of different frequency (MHz) domain in Watts:
Frequency     total_pwr leakage_pwr internal_pwr switching_pwr %_total_pwr
4.000e+02     1.729e-02 1.704e-04    1.107e-02     6.052e-03    8.333e+01%
2.000e+02     3.4595e-03 2.1869e-05 3.1000e-03     3.3758e-04   1.666e+01%
0.000e+00     0.000e+00 0.000e+00    0.000e+00     0.000e+00    0.000e+00%

Power of different Vdd domain in Watts:
Vdd_domain total_pwr leakage_pwr internal_pwr switching_pwr %_total_pwr
VDD (1.1V) 2.075e-02 1.922e-04 1.417e-02      6.389e-03      1.000e+02%

Power of different cell types in Watts:
cell_type    total_pwr leakage_pwr internal_pwr switching_pwr %_total_pwr
combinational 7.111e-03 1.221e-04 1.754e-03       5.235e-03 3.426e+01
latch_and_FF 4.292e-03 3.890e-05 3.323e-03        9.301e-04 2.067e+01
memory        0.000e+00 0.000e+00 0.000e+00       0.000e+00 0.000e+00
I/O           0.000e+00 0.000e+00 0.000e+00       0.000e+00 0.000e+00
clocked_inst 9.352e-03 3.119e-05 9.096e-03        2.246e-04 4.505e+01
decap         0.000e+00 0.000e+00 0.000e+00       0.000e+00 0.000e+00

where clocked_inst are instances that cannot be classified as
latch_and_FF, memory, or I/O, but have clock pin(s).

Total chip power, 0.020757 Watt including core power and other
domain power.
Total clock network only power, 0.003225 Watt. Total clock power,
including clock network and FF/latch clock pin power, 0.005992 Watt.
\paragraph{power\_summary.rpt 字段说明}
\begin{itemize}
  \item \texttt{Recommended dynamic simulation time}：建议的 \gsr{DYNAMIC_SIMULATION_TIME}，以覆盖大部分功耗事件
  \item 按频率域、Vdd 域、单元类型列出 \texttt{total\_pwr}、\texttt{leakage\_pwr}、\texttt{internal\_pwr}、\texttt{switching\_pwr} 及占总功耗百分比
  \item \texttt{clocked\_inst}：有 clock pin 但无法归类为 latch\_and\_FF、memory 或 I/O 的实例
  \item \texttt{Total chip power}：含核与其它域的总功耗
  \item \texttt{Total clock network only power} / \texttt{Total clock power}：仅时钟网络 / 含 FF-latch clock pin 的时钟功耗
\end{itemize}

\end{lstlisting}

\subsection{\texttt{<top\_level\_block>.power.rpt} 文件}
\subsubsection*{\texttt{<top\_level\_block>.power.rpt} File}

功耗计算在 \texttt{adsRpt} 目录创建 \texttt{<top\_level\_block>.power.rpt}，包含每个 instance 的功耗信息（功耗单位为 W，电流为 A）。多 Vdd/Vss 设计中，每个电源 pin 有独立数据组（末尾为 pin 名）：
\begin{lstlisting}
<inst_name> <cell_name> <freq_Hz> <toggle_rate> <leakage_power>
<switching_power> <internal_power + clock_pin_power> <total_power>
<X_loc_um> <Y_loc_um> <VDD/VSS_domain>
<[leakage_data_source]_[internal_pwr_data source]>
<domain_type: [ 1 | 2 | 0 ]> <leakage_current> <total_current>
<LIB_name> <cell_p/g_pin_name>
\end{lstlisting}
其中：
\begin{itemize}
  \item \texttt{<inst\_name>}：instance 名
  \item \texttt{<cell\_name>}：单元名
  \item \texttt{<freq\_Hz>}：工作频率
  \item \texttt{<toggle\_rate>}：每时钟周期平均翻转次数
  \item \texttt{<leakage\_power>}：漏电功耗
  \item \texttt{<switching\_power>}：开关功耗
  \item \texttt{<internal\_power + clock\_pin\_power>}：内部功耗（含时钟 pin）
  \item \texttt{<total\_power>}：该 pin 总功耗
  \item \texttt{<X\_loc\_um> <Y\_loc\_um>}：坐标
  \item \texttt{<VDD/VSS\_domain>}：电源域名
  \item \texttt{<[leakage\_data\_source]\_[internal\_pwr\_data source]}：各为 \texttt{apl} 或 \texttt{lib}
  \item \texttt{<domain\_type>}：1 = 非 multi\_rail\_pwr\_domain；2 = multi\_rail\_Vdd\_domain；0 = multi\_rail\_VSS\_domain
  \item \texttt{<leakage\_current>}：漏电电流
  \item \texttt{<total\_current>}：总电流（A）
  \item \texttt{<LIB\_name>}：单元 .lib 名（未在 .lib 中定义则显示 \texttt{lef2lib}）
  \item \texttt{<cell\_pg\_pin\_name>}：P/G pin 名
\end{itemize}

\subsection{\texttt{apache.power.info} 文件}
\subsubsection*{\texttt{apache.power.info} File}

\texttt{adsRpt} 目录下创建 \texttt{apache.power.info}，包含 VCD 与 STA 文件中缺失网络等附加信息，用于调试功耗计算结果。常见内容包括：未在 VCD/STA 中出现的网络列表、功耗缩放警告摘要、以及实例级数据不一致提示。

\paragraph{apache.power.info 典型用途}
\begin{itemize}
  \item 核对 VCD 与 DEF 名称映射是否完整
  \item 查找 STA 未覆盖但功耗显著的实例
  \item 确认 \gsr{INSTANCE_POWER_FILE} 与设计中实例名是否一致
\end{itemize}

\section{静态与动态电压、电流分析结果}
\subsection*{Results for Static and Dynamic Voltage and Current Analyses}

\subsection{EM、静态 IR 与 DvD 色阶图显示}
\subsubsection*{EM, Static IR and DvD Colormap Displays}

多数 RedHawk 分析结果可在 GUI 中通过相应菜单命令或按钮显示为全芯片色阶图（见附录 D GUI 说明），或使用 TCL \cmd{show} 命令（见附录 D「TCL / Script Commands」）。

\subsection{静态 EM 与 IR Drop 结果文件}
\subsubsection*{Static EM and IR Drop Results Files}

以下静态摘要文件位于 \texttt{adsRpt/Static} 目录。
\begin{noteBox}
所列网络名为实际域名。
\end{noteBox}

\texttt{<design>.em.worst} 按恶化程度递减列出导线段与 via 的最差 EM 违例。每行格式：
\begin{itemize}
  \item 导线：\texttt{<metal\_layer> <location> <EM\_ratio> <domain\_name> <metal\_width>}
  \item Via：\texttt{<via\_name> <location> <EM\_ratio> <VDD | VSS>}
\end{itemize}

\texttt{<design>.ir.worst} 按恶化程度递减列出静态 IR drop（电压跌落与地弹）最严重的节点。每行：
\texttt{<actual\_voltage> <ideal\_voltage> <domain\_name> <x\_y\_location> <layer\_name>}。

示例：
\begin{lstlisting}
#voltage #volt #net           #x_y_location       #layer_name
1.0734     1.0800 VDD2 ( 261.700, 155.760)MET2
1.0734     1.0800 VDD2 ( 262.780, 155.760)MET2
1.0736     1.0800 VDD2 ( 261.700, 232.980)MET2
...
\end{lstlisting}

\texttt{<design>.inst.worst} 列出静态 IR drop 最严重的 instance。每行：
\texttt{<inst\_vdd> <vdd\_drop> <gnd\_bounce> <ideal\_vdd> <domain\_name> <x\_y\_loc> <inst\_name>}。

示例：
\begin{lstlisting}
#inst_vdd #inst_drop #gnd_bounce #ideal_vdd #pwr_net #x_y_location #inst_name
1.0787 0.0006 0.0006 1.0800 VDD1( 60.7200, 273.7350) inst1
1.0787 0.0006 0.0006 1.0800 VDD1( 71.6100, 273.7350) inst2
1.0786 0.0008 0.0006 1.0800 VDD1( 28.3800, 265.1550) inst3
...
\end{lstlisting}

\texttt{switch\_static.rpt} 仅用于低功耗设计中的 header/footer 开关 instance，给出开关的电压与电流信息。每行：
\texttt{<internal\_node\_voltage> <voltage\_on\_switch> <average\_current> <header | footer> <instance\_name>}。

示例：
\begin{lstlisting}
#int_node_volt #volt_on_switch #avg_current #type #instance_name
1.0800 0.0037        0.6248 header QNP_1_CORNER_HS_CELL_002
1.0800 0.0033        0.5533 header QNP_1_CORNER_HS_CELL_001
1.0796 0.0049        0.8231 header QNP_1_BOTTOM_HS_CELL_115
...
\end{lstlisting}

\subsection{静态与动态电压跌落结果文件}
\subsubsection*{Static and Dynamic Voltage Drop Results Files}

以下动态摘要文件位于 \texttt{adsRpt/Dynamic/} 目录。长时间运行中可用 \cmd{xgraph <filename>} 查看 ivdd、ipwr 等中间图形结果。
\begin{noteBox}
所列网络名为实际域名。
\end{noteBox}

\paragraph{\texttt{<design>.ir.worst} 文件}
列出最差动态电压跌落位置，每行格式与静态 \texttt{.ir.worst} 相同。示例：
\begin{lstlisting}
#voltage #ideal_volt              #net #x_y_location    #layer_name
0.9262           1.0800           VDD1 (261.700, 155.760)   MET2
0.9275           1.0800           VDD1 (262.780, 155.760)   MET2
0.9504           1.0800           VDD1 (174.500, 155.760)   MET2
...
\end{lstlisting}

\paragraph{\texttt{<design>.dvd} 文件}
按 VDD-VSS 差分报告各 instance 的电压跌落，按恶化程度递减。每行：
\begin{lstlisting}
<x_loc> <y_loc> <effective_vdd-vss_over_TW> <max_vdd-vss_over_TW>
<min_vdd-vss_over_TW> <min_vdd-vss_of_clockcycle>
<min_vdd_voltage_over_TW> <max_gnd_bounce_over_TW> <VDD | VSS>
<instance_name>
\end{lstlisting}
为获得有效 Vdd，数值在若干小采样窗口上平均，而非整个 STA 时序窗口。通过在整个时序窗口上滑动小窗口得到多个平均有效 Vdd，报告其中最差平均值作为有效 (Vdd-Vss)。若无时序窗口数据，则不报告基于时序窗口的 DvD 值。

\texttt{<TW\_flag>} 表示动态仿真时间对时序窗口（Timing Window）的覆盖情况：
\begin{itemize}
  \item 空字符串：TW 在动态仿真时间内被完全覆盖
  \item \texttt{\^{}}：TW 被部分覆盖
  \item \texttt{*}：TW 在动态仿真时间内未被覆盖
\end{itemize}
各列含义：\texttt{eff\_vdd} 为时序窗口内有效 (Vdd-Vss) 平均值（最差）；\texttt{max\_pg\_tw}/\texttt{min\_pg\_tw} 为 TW 内最大/最小 Vdd-Vss；\texttt{min\_pg\_sim} 为整仿真周期最小 Vdd-Vss；\texttt{min\_vdd\_tw}/\texttt{max\_vss\_tw} 为 TW 内最小 Vdd 与最大地弹。

示例：
\begin{lstlisting}
#loc_x loc_y eff_vdd max_pg_tw min_pg_tw min_pg_sim min_vdd_tw max_vss_tw
domain instance_name
218.122 167.320 0.9671 0.9794 0.9555 0.8798 1.0012 0.0458 VDD inst241
267.134 157.905 0.9808 0.9904 0.9736 0.9181 1.0122 0.0386 VDD inst3469
264.996 153.615 1.0034 1.0079 0.9992 0.9635 1.0069 0.0085 VDD inst2463
...
\end{lstlisting}

\subsection{Via 的静态与动态结果}
\subsubsection*{Static and Dynamic Results for Vias}

RedHawk 可为静态与动态分析生成 Via 电压跌落与电流报告，位于 \texttt{adsRpt} 文件夹。可通过 GSR 关键字 \gsr{VIA_IR_REPORT 1} 开启，或使用 TCL：
\begin{lstlisting}
report [ ir | dvd ] -via -o <output_file>
\end{lstlisting}
输出文件名：
\begin{itemize}
  \item 静态：\texttt{adsRpt/Static/<>.via.ir.worst}
  \item 动态：\texttt{adsRpt/Dynamic/<>.via.ir.worst}
\end{itemize}

\figplaceholder{Figure 6-10}{Via 电压跌落与电流报告}

\paragraph{Via IR 报告字段}
静态/动态 via 报告每行含 via 名、位置坐标、IR/EM 比值、所属 VDD 或 VSS 域等，用于定位过孔处的压降与电流违例热点，与线网 \texttt{.ir.worst} 互补。

\subsection{电流报告文件}
\subsubsection*{Current Report Files}

\paragraph{\texttt{<design>.ignd}}
汇总仿真时间内所有 instance 的 Vss 电流。每行：\texttt{<simulation\_time> <current\_value>}。示例：
\begin{lstlisting}
Time            i(gnd)
0.000000        0.0
-3480           -0.0038748
-3470           -0.00309187
-3460           -0.00320901
...
\end{lstlisting}

\paragraph{\texttt{<design>.ipwr}}
汇总所有 instance 的 Vdd 电流需求。每行：\texttt{<simulation\_time> <current\_value>}。示例：
\begin{lstlisting}
Time           i(pwr)
0.000000       0.0
-3480          0.0038748
-3470          0.00309187
...
\end{lstlisting}

\paragraph{\texttt{<design>.ivdd}}
报告芯片电源电流。该电流由各电源/地源的电压源提供。无封装时，ivdd 测量设计 pad 处电流，包含片上寄生、decap 及其他电容效应。有封装时，ivdd 测量封装外电压源处电流，另含封装寄生。
\paragraph{\texttt{<design>.ivdd} 示例}
无封装时测量设计 pad 电流（含片上寄生、decap 等）；有封装时测量封装外电压源电流（另含封装寄生）：
\begin{lstlisting}
Time                 i(vdd)
0.000000           0.0
-3490.0000         0.00265285
-3480.0000         0.0036665
-3470.0000         0.00341904
...
\end{lstlisting}


\paragraph{\texttt{<design>.ivdd.vsrc}}
动态分析后在 \texttt{adsRpt/Dynamic} 生成，包含所有电源源的供电电流波形（图~\ref{fig:rh-current-waveforms}）。绘制指定网络所有 pad 总供电电流的 TCL 命令：
\begin{lstlisting}
plot current -net -name <net_name> -pad ?-sv? ?-xgr?
\end{lstlisting}

\figplaceholder{Figure 6-11}{电流需求与供电电流波形}
\label{fig:rh-current-waveforms}

\paragraph{\texttt{<design>.ipwr.domain}}
动态分析后生成，包含 GSR \gsr{VDD_NETS} 定义的各电源域电流需求波形。

\paragraph{\texttt{<design>.ignd.domain}}
动态分析后生成，包含 GSR \gsr{VSS_NETS} 定义的各地电源域电流波形。

\paragraph{\texttt{switch\_dynamic.rpt}}
仅用于低功耗 header/footer 开关。矢量无关开关分析列出最差 VDD-VSS 差分及仿真时间内开关电压、电流最大值。每行：instance 名、开关类型、最差 Vdd-Vss 电压、最大开关电压与电流。

Ramp-up 分析报告包含最大开关电流及仿真结束时各 header/footer 开关状态（OFF 或 ON）：
\begin{lstlisting}
#instance_name type maximal_Isw(Amp) status (SAT)
#SAT: maximal_Isw > saturation_current as specified in switch model file
switch_inst_R17_C52 header 5.000000e-09 OFF SAT
switch_inst_R17_C52 header 2.603736e-09 OFF
\end{lstlisting}
可快速查看哪些开关已 ON 及其效率。
\paragraph{\texttt{switch\_dynamic.rpt} 示例（vectorless）}
每行：实例名、开关类型、最差 Vdd-Vss、最大开关电压与电流：
\begin{lstlisting}
#<inst_name> <type> <worst_int_Vdd-Vss_Volt> <max_Vsw_Volt> <max_Isw_Amp>
QNP_1_BOTTOM_HS_CELL_083 H 1.051716 0.024256 0.004053
QNP_1_BOTTOM_HS_CELL_082 H 1.051716 0.025019 0.004180
...
\end{lstlisting}


\paragraph{\texttt{decaps.rpt}}
取决于 GSR \gsr{DYNAMIC_REPORT_DECAP}：值为 1 时报告 GSR 定义的全部 intentional decap 的最大电流；值为 2 时另报告非开关 instance 的最大电流；值为 0（默认）时不生成 decap 报告。

\paragraph{\texttt{freqd\_ipwr.out}}
baseline DvD 流程生成；其他模式不生成。各 instance 可关联特定时钟与频率域。该文件为各频率域内所有开关电流之和的波形，可用：
\begin{lstlisting}
xgraph freqd_ipwr.out
\end{lstlisting}
显示。文件两列：时间与电流（A）。

\subsection{Pad 电流文件}
\paragraph{静态 \texttt{pad.current} 示例}
\begin{lstlisting}
#current      #pad_center_location                  #pad_name
2.4745        ( 65.282, 0.082)                  Vdd_130564_1645
2.1274        ( 163.329, 0.082)                 Vdd_326658_1645
...
\end{lstlisting}

\paragraph{动态 \texttt{pad.current} 示例}
\begin{lstlisting}
Vss1
-1750 -1.13646e-05
-1740 -0.00240712
...
0   -0.0506185
Vdd1
-1750 1.13646e-05
...
\end{lstlisting}

\subsubsection*{Pad Current File}

\texttt{pad.current} 指定 pad 单元或 pad 位置的电流（不含 presim 时间电流），用于判断指定 pad 位置连通是否完整。静态 IR/EM 与动态压降分析均生成该报告。静态报告写入 \texttt{adsRpt/Static}，每行：
\texttt{<current\_value> <x\_y\_location\_of\_pad\_center> <pad\_name>}。

动态报告写入 \texttt{adsRpt/Dynamic}，格式为各 pad 在仿真时间上的电流：
\begin{lstlisting}
<pad_name_a>
<simulation_time1> <current_value>
...
<pad_name_b>
<simulation_time1> <current_value>
...
\end{lstlisting}

\subsection{\texttt{probe\_nodes.out} 文件}
\subsubsection*{\texttt{probe\_nodes.out} file}

\texttt{probe\_nodes.out} 报告静态与动态分析的探针节点电压与电流，分别位于 \texttt{adsRpt/Static/} 或 \texttt{adsRpt/Dynamic/}。须在 \gsr{PROBE_NODE_FILE} GSR 关键字指定文件中定义探针位置。示例：
\begin{lstlisting}
#Probe_Name                 I                 V
cell_198_VDD_INT    3.454429e-05           1.091716e+00
cell_198_VSS        2.562570e-05           9.402982e-02
\end{lstlisting}

\section{CMM 约束违例报告}
\subsection*{CMM Constraint Violation Reports}

\subsection{约束违例摘要}
\subsubsection*{Constraint Violation Summary}

顶层分析后，若子 block CMM 中定义了约束，仿真结果后处理会在 \texttt{redhawk.log} 中记录违例摘要表，例如：
\begin{lstlisting}
Summary of CMM Constraint Violation:
---------------------------------------------------
Constraint Type No. of violations
---------------------------------------------------
Top Connection Min Voltage 50
Top Connection Max Voltage 50
Pin Min Voltage 500
Pin Max Voltage 500
There are total 1100 constraint violations in the CMM
instances, please look at adsRpt/Dynamic/
CMM_constraint_violation.rpt for detail!
\end{lstlisting}

\subsection{动态分析约束摘要}
\subsubsection*{Dynamic Analysis Constraint Summary}

动态分析的最终 CMM 约束违例报告位于 \texttt{adsRpt/Dynamic/CMM\_constraint\_violation.rpt}。格式包括：
\begin{itemize}
  \item Top connection max/min voltage constraint violation
  \item pin max/min voltage constraint violation
\end{itemize}
每类违例注释行说明：\texttt{<inst> <cell> <x,y,layer> <net\_name> <constraint> <voltage\_value>} 或 \texttt{<inst:pin> <cell> <x,y> <constraint> <voltage\_value>}。

\subsection{静态分析约束摘要}
\subsubsection*{Static Analysis Constraint Summary}

静态分析最终报告位于 \texttt{adsRpt/Static/CMM\_constraint\_violation.rpt}，格式为 Top connection IR constraint violation 与 pin IR constraint violation
\paragraph{动态 CMM 违例报告示例}
\begin{lstlisting}
# Top connection min voltage constraint violation:
ABC_CORE/ram64/adsU1 ram32x8x2 (1759.3600,783.0350, METAL2) VDD 1.7820 1.7309
# Top connection max voltage constraint violation:
ABC_CORE/ram64/adsU1 ram32x8x2 (1763.9600,778.4350, METAL2) VSS 0.0010 0.0528
# Pin min voltage constraint violation:
ABC_CORE/ram64/adsU1:VDD ram32x8x2 (1792.3700,727.8150) 1.7460 1.7294
\end{lstlisting}

\paragraph{静态 CMM 违例报告示例}
\begin{lstlisting}
# Top connection IR constraint violation:
ABC_CORE/ram64/adsU1 ram32x8x2 (1759.3600,783.0350, METAL2) VDD 1.7960 1.7909
# Pin IR constraint violation:
ABC_CORE/ram64/adsU1:VDD ram32x8x2 (1792.3700,727.8150) 1.7910 1.7909
\end{lstlisting}
。

\section{在 GUI 中使用摘要文件调试}
\subsection*{Debugging Using Summary Files in the GUI}

生成文本结果报告后，可在 GUI 中交互调试。以下 GUI 命令可生成违例排序列表，选中单项可在版图上高亮：
\begin{itemize}
  \item Results $>$ List of worst EM for static simulation
  \item Results $>$ List of worst IR instances for static simulation
  \item Results $>$ List of worst IR for wire and via for static simulation
  \item Results $>$ List of worst DvD instances for dynamic simulation
\end{itemize}

\section{多 Vdd/Vss 分析的输出文件}
\subsection*{Output Files from Multiple Vdd/Vss Analysis}

表~\ref{tab:multi-vdd-files} 汇总多 Vdd 分析可用文件与报告。

\begin{table}[htbp]
\centering
\caption{多 Vdd 结果文件}
\label{tab:multi-vdd-files}
\small
\begin{tabular}{@{}p{5.5cm}p{7.5cm}@{}}
\toprule
文件名 & 信息 \\
\midrule
\texttt{adsRpt/apache.refCell.noPGarcs} & 列出缺少 P/G arc 信息的单元及 P/G pin 名 \\
\texttt{adsRpt/Dynamic/<design>.dvd.arc} & 与标准 \texttt{.dvd} 相同，但按 arc 报告 \\
\texttt{adsRpt/Dynamic/<design>.dvd.pin} & 与标准 \texttt{.dvd} 相同，但按 pin 报告 \\
\texttt{adsRpt/Static/<design>.inst.pin} & 静态 instance 压降，每个 instance 每个 P/G pin 一行 \\
\texttt{adsRpt/Static/<design>.inst.arc} & 静态 instance 压降，每个 instance 每个 P/G arc 一行 \\
\bottomrule
\end{tabular}
\end{table}

\paragraph{Explorer 与 GUI 调试菜单}
除第 6 章所列文本报告外，\textbf{Explorer} 菜单提供更完整的结果浏览与后处理；\textbf{Results} 菜单可生成最差 EM/IR/DvD 排序列表并在版图中高亮。建议签核前同时检查 \texttt{adsRpt/Error}、\texttt{adsRpt/Warn} 与 \texttt{power\_summary.rpt}。

\section{其他文件}
\subsection*{Other Files}

RedHawk 在 \texttt{adsRpt} 目录写出多种文件，为分析运行前后提供附加信息。

\subsection{杂项}
\subsubsection*{Miscellaneous}

以下文件报告非标准库或仿真条件：

\begin{table}[htbp]
\centering
\small
\begin{tabular}{@{}p{4.5cm}p{8cm}@{}}
\toprule
文件 & 说明 \\
\midrule
\texttt{apache.noLefPins} & LEF macro 中未定义 pin 的单元 \\
\texttt{apache.noLibPins} & 库定义中无 pin 的单元 \\
\texttt{apache.rcNoDriver} & 无驱动器的网络 \\
\texttt{apache.refCell.noLefLib} & 设计中存在但不在库或 LEF 中的单元 \\
\texttt{apache.refCell.noLib} & 设计中存在但不在库中的单元 \\
\texttt{apache.refCell.noLef} & 设计中存在但不在 LEF 中的单元 \\
\texttt{apache.refCell.noPwr} & 库中无内部功耗表的单元 \\
\texttt{apache.refCell.noTriggerEdge} & 库中无可识别时钟触发沿的单元 \\
\texttt{apache.inst.libCurrent} & 缺 APL 数据的 instance 的 .lib 电流剖面 \\
\texttt{apache.switchCellnoTW} & 开关单元但 STA 中控制 pin 无时序窗口 \\
\texttt{clock.untraced} & STA 时钟网络但 RedHawk 无法追踪的网络 \\
\texttt{<design>.power.unused} & 未接入电源网络的 instance \\
\texttt{<design>.power.worst} & 静态功耗最高的 instance \\
\texttt{<design\_name>.rcFail} & 无法为输出负载创建 pi-model 的网络 \\
\texttt{<design\_net>.Via.unconnect} & 与该网络所属、与电源物理隔离的 via \\
\texttt{<design\_net>.Wire.unconnect} & 与驱动源物理隔离的导线 \\
\texttt{<design\_net>.Pin.unconnect} & 与电源物理隔离的 pin \\
\texttt{<design\_net>.PinInst.unconnect} & 与电源物理隔离的 pin instance \\
\texttt{<design\_net>.PinInst.unconnect.logic} & 逻辑断开但物理连通的 PinInst 及所连网络 \\
\texttt{<design>.unplaced} & DEF 中标记 UNPLACED 的 instance \\
\texttt{<design>.stackvia} & 工作目录或 adsRpt 中缺失的叠层 via \\
\texttt{GC.unreachable.fflatch} & 门控时钟控制不可达的 FF/latch 比例与列表 \\
\texttt{libParser.errMsgs} & 解析 Synopsys Liberty 库时的错误 \\
\texttt{ppi\_multi\_group\_pin\_geo\_cells} & pin 有多组断开几何且 IR drop 显著的单元 \\
\texttt{undefined\_cells} & 设计中实例化但 DEF/LEF 无定义的 master cell \\
\texttt{PG.ploc} & 提取时识别的 P/G 源、坐标与层；使用 .pad/.pcell 时，若 pad 与布线同层有多处连接，源名格式为 \texttt{<instance>:<net>:<\#>}；导入 *.ploc 时内容应一致，除非存在断连源 \\
\texttt{DEF.ploc} & 顶层 DEF PINS 段的重格式化版本；\gsr{ADD_PLOC_FROM_TOP_DEF} 为 1 时自顶层 DEF 读取源 \\
\bottomrule
\end{tabular}
\end{table}

\begin{seeAlsoBox}
若脚本同时包含 import pad 命令与 \gsr{ADD_PLOC_FROM_TOP_DEF 1}，RedHawk 将退出。
\end{seeAlsoBox}

\subsection{调试}
\subsubsection*{Debugging}

\begin{table}[htbp]
\centering
\small
\begin{tabular}{@{}p{4.5cm}p{8cm}@{}}
\toprule
文件 & 说明 \\
\midrule
\texttt{apache.clkPin0} & 连接到非时钟网络的时钟 pin \\
\texttt{apache.dupLIBCells} & 库中多次引用的单元 \\
\texttt{apache.rc0Net} & 电容为零的网络 \\
\texttt{apache.rcBogus} & SPEF 中无法映射到设计的网络与行 \\
\texttt{apache.rcMismatch} & SPEF 与 DEF 驱动不一致报告 \\
\texttt{apache.slew0} & 时序文件中缺 slew 的 instance；时序单元列出 CLK/OUT pin，组合单元列出 IN/OUT pin \\
\texttt{apache.staBogus} & 错误 STA 数据（如设计中不存在的 instance 的 TW） \\
\texttt{apache.staFlatten} & RedHawk 库内被展平且 TW 被忽略的 instance \\
\texttt{apache.tw0} & STA 文件未覆盖的 instance \\
\texttt{apache.twclk0} & 时钟 pin 无时序窗口覆盖的时序单元 \\
\texttt{apache.twclkLate} & 输出 pin 在时钟 pin 之前翻转的时序单元 \\
\texttt{inst.reactivated} & STA 未覆盖但 toggle rate 被强制非零的 instance \\
\texttt{<design>\_<net>.collided\_vias} & 与其他 via 接触或碰撞的 via \\
\texttt{aplchk.log} & APL 相关错误与警告 \\
\bottomrule
\end{tabular}
\end{table}

\subsection{动态仿真准备}
\subsubsection*{Dynamic Simulation Preparation}

\begin{table}[htbp]
\centering
\small
\begin{tabular}{@{}p{4.5cm}p{8cm}@{}}
\toprule
文件 & 说明 \\
\midrule
\texttt{Dynamic/cell.rpt} & 有/无 APL 表征数据的单元列表 \\
\texttt{<design>.PG.unconnect} & 与 VDD、GND 均物理隔离的 instance（动态分析忽略） \\
\texttt{<design>.VDD.unconnect} & 与 VDD 物理隔离的 instance \\
\texttt{<design>.GND.unconnect} & 与 GND 物理隔离的 instance \\
\texttt{<design>.noTiming} & 无有效时序窗口的 instance（动态分析不翻转） \\
\bottomrule
\end{tabular}
\end{table}

\subsection{低功耗 Ramp Up 分析}
\subsubsection*{Low Power Ramp Up Analysis}

\begin{table}[htbp]
\centering
\small
\begin{tabular}{@{}p{4.5cm}p{8cm}@{}}
\toprule
文件 & 说明 \\
\midrule
\texttt{virtual\_domain\_total\_i.rpt} & 各由 header/footer 开关控制的电压域总 ramp-up 电流波形（时间、电流 A） \\
\texttt{virtual\_domain\_worst\_v.rpt} & 各电压域 ramp-up 期间最差节点电压波形（时间、电压 V） \\
\bottomrule
\end{tabular}
\end{table}
